% generated by roam.py handbook export -- не редактировать, править заметки
% 20260915183320-собственные_значения_симметричных_матриц.org (body b0b3509b735f)

\begin{theorem}[{Собственные значения симметричных матриц}]\label[theorem]{note:4F4E5575-ABE6-4D9B-9C61-B76A3AA9701C}
  У вещественных симметричных матриц все \hyperref[note:9D2905FC-8E68-4E6F-9243-364269F795B0]{собственные значения}~--- вещественны
\end{theorem}

\begin{proof}
  Пусть $A = A^T$~--- вещественная матрица, $Aq = \lambda q$, $q \ne 0$. Здесь $\lambda \in \C$ и $q \in \C^n$.
  Обозначим $q^* = \overline{q}^T$ и посчитаем число $q^* A q$:
  \[
  q^* A q = q^* (\lambda q) = \lambda q^* q = \lambda \norm{q}^2
  \]
  Это число совпадает со своим сопряженным:
  \[
  \overline{q^* A q} = (q^* A q)^* = q^* A^* q = q^* A q
  \]
  так как $A^* = \overline{A}^T = A^T = A$ ($A$~--- вещественная и симметричная). Значит $q^* A q \in \R$.

  Так как $\norm{q}^2 > 0$~--- вещественное число, то
  \[
  \lambda = \frac{q^* A q}{\norm{q}^2} \in \R
  \]
\end{proof}
